% !TeX root = part7.tex
\documentclass[../final.tex]{subfiles}

\begin{document}

% Конспект: Desktop/4.jpg, 5.jpg (17.10) и низ 2.jpg (формула для <q^2>).
% В тетради g_2D = S w/(pi v0^2): потерян множитель 1/2, исправлено.
% Расходимости в 1D и 2D уточнены: какая часть (нулевая или тепловая)
% и как именно растёт <q^2> с размером.
\practice{7}{17.10.2026}{Средний квадрат смещения. Плотность состояний}
\captionsetup{font=footnotesize,labelfont=bf,skip=5pt}

\begingroup
\renewcommand{\thesection}{4.\arabic{section}}
\task[17]{Средний квадрат смещения атома в решётке}
\label{task:fksv_displacement}

\condition
Определить средний квадрат смещения атома в решётке.
Решить ту же задачу, используя теорему вириала.
Продемонстрировать расходимость в 1D и 2D.

\solution
Полагаем \rassumption{$k_{\mathrm B}=1$} и, как на прошлом занятии,
фиксируем центр масс: мода $k=0$ исключена.

\subsection*{Через операторы рождения и уничтожения}

Смещение атома $r$ выражается через нормальные координаты:
\[
    \hat q_r=\frac1{\sqrt N}\sum_k\hat Q_ke^{ikr},\qquad
    \hat Q_k=\sqrt{\frac{\hbar}{2m\omega_k}}
       \left(\hat b_k+\hat b_{-k}^\dagger\right).
\]
Средние $\langle\hat Q_k\hat Q_{k'}\rangle$ отличны от нуля только
при $k'=-k$, и тогда из тепловых средних прошлого занятия
\[
    \langle\hat Q_k\hat Q_{-k}\rangle
       =\frac{\hbar}{2m\omega_k}
        \langle\hat b_k\hat b_k^\dagger+\hat b_{-k}^\dagger\hat b_{-k}\rangle
       =\frac{\hbar}{m\omega_k}\left(\bar n_k+\frac12\right).
\]
Следовательно,
\[
    \langle q^2\rangle=\frac1N\sum_k\langle\hat Q_k\hat Q_{-k}\rangle
       =\sum_k\frac{\hbar}{Nm\omega_k}
         \left(\frac12+\frac1{e^{\hbar\omega_k/T}-1}\right).
\]
По трансляционной симметрии ответ не зависит от номера атома.

\subsection*{Через теорему вириала}

Для одной частоты $\omega$ энергия $N$ осцилляторов
\[
    E=N\hbar\omega\left(\frac12+\frac1{e^{\hbar\omega/T}-1}\right).
\]
Для гармонического осциллятора средние кинетическая и потенциальная
энергии равны, поэтому
\[
    \langle U\rangle=\frac E2
       =\frac12N\hbar\omega\left(\frac12+\bar n\right).
\]
С другой стороны,
$\langle U\rangle=\sum_r\frac12m\omega^2\langle q_r^2\rangle
=\frac12Nm\omega^2\langle q^2\rangle$, откуда
\[
    \langle q^2\rangle=\frac{\hbar}{m\omega}\left(\frac12+\bar n\right).
\]
В решётке каждая мода $k$ даёт то же соотношение для
$\langle\hat Q_k\hat Q_{-k}\rangle$, а
$\sum_r\hat q_r^2=\sum_k\hat Q_k\hat Q_{-k}$ (равенство Парсеваля).
Деля на $N$, получаем ту же сумму по $k$, что и выше:
вклад моды равен $\hbar(\bar n_k+1/2)/(Nm\omega_k)$.

\subsection*{Переход к интегралу}

Перейдём к размерному волновому вектору $\boldsymbol k$ и рассмотрим
\rassumption{одну акустическую ветвь $\omega_k=ck$}: расходимости
определяются длинными волнами. Число состояний в элементе
фазового объёма
\[
    d\mathcal N=\frac{d^dx\,d^dp}{(2\pi\hbar)^d}
       =\frac{L^d\,d^dk}{(2\pi)^d},\qquad
    \boldsymbol p=\hbar\boldsymbol k,\qquad N=\left(\frac La\right)^d.
\]
В образце размера $L$ наименьшее волновое число
$k_{\min}\sim2\pi/L$; при $L\to\infty$ интегралы идут от нуля.

\subsubsection*{Одномерная цепочка}
\[
    \langle q^2\rangle
       =2\int_{k_{\min}}^{\pi/a}\frac{L\,dk}{2\pi}\,
          \frac{\hbar}{Nmck}
          \left(\frac12+\frac1{e^{\hbar ck/T}-1}\right).
\]
Интеграл \rnote{расходится на нижнем пределе}. Уже нулевые колебания
дают $\int dk/k$:
\[
    \langle q^2\rangle_{T=0}\simeq
       \frac{\hbar a}{2\pi mc}\ln\frac{L}{a}.
\]
При $T>0$ для длинных волн $\bar n_k\simeq T/(\hbar ck)\gg1$,
подынтегральное выражение ведёт себя как $T/k^2$, и
$\langle q^2\rangle$ растёт пропорционально $L$. Точная сумма для
классической цепочки ($\omega_j^2=(4\alpha/m)\sin^2(\pi j/N)$):
\[
    \langle q^2\rangle=\frac{T}{mN}\sum_{j=1}^{N-1}\frac1{\omega_j^2}
       =\frac{T}{4\alpha N}\sum_{j=1}^{N-1}\frac1{\sin^2(\pi j/N)}
       =\frac{T}{12\alpha}\left(N-\frac1N\right).
\]

\subsubsection*{Двумерная решётка}
\[
    \langle q^2\rangle
       =\int_{k_{\min}}\frac{L^2\,2\pi k\,dk}{(2\pi)^2}\,
          \frac{\hbar}{Nmck}
          \left(\frac12+\frac1{e^{\hbar ck/T}-1}\right).
\]
Нулевые колебания здесь дают сходящийся интеграл $\int dk$.
Тепловая часть ведёт себя как $\int dk/k$ и
\rnote{логарифмически расходится}:
\[
    \langle q^2\rangle_T\simeq\frac{a^2T}{2\pi mc^2}\ln\frac{L}{a}.
\]

\subsubsection*{Трёхмерный кристалл}
Множитель $k^2\,dk$ делает оба интеграла сходящимися:
смещение конечно при любой температуре.

\answer
\[
    \langle q^2\rangle=\sum_k\frac{\hbar}{Nm\omega_k}
       \left(\frac12+\frac1{e^{\hbar\omega_k/T}-1}\right).
\]
\rresult{В 1D смещение расходится уже при $T=0$ (как $\ln L$),
а при $T>0$ растёт как $L$; в 2D при $T>0$ растёт как $\ln L$;
в 3D конечно.} Тепловые длинноволновые колебания разрушают дальний
порядок в одномерных и двумерных кристаллах.
\endgroup

\rtheory{Плотность состояний}

Суммы по модам удобно переписать через \rconcept{плотность состояний}
\[
    g(\omega)=\sum_k\delta(\omega-\omega_k),\qquad
    g(\omega)\,d\omega=d\mathcal N\ \text{--- число мод в }(\omega,\omega+d\omega).
\]
Тогда тепловая энергия решётки
\[
    E=\int\frac{\hbar\omega\,g(\omega)\,d\omega}{e^{\hbar\omega/T}-1}.
\]

\subsection*{Линейный закон дисперсии}

Пусть $\omega_k=v_0k$. В одном измерении $k$ пробегает оба знака:
\[
    g_{1\mathrm D}(\omega)=\frac L{2\pi}\int_{-\infty}^{\infty}
       \delta(\omega-v_0|k|)\,dk
     =\frac L\pi\int_0^\infty\frac1{v_0}
       \delta\left(k-\frac\omega{v_0}\right)dk
     =\frac L{\pi v_0}.
\]
В двух измерениях, на одну ветвь,
\[
    g_{2\mathrm D}(\omega)=\frac S{(2\pi)^2}\int_0^\infty
       \delta(\omega-v_0k)\,2\pi k\,dk
     =\frac S{2\pi v_0}\cdot\frac\omega{v_0}
     =\frac{S\omega}{2\pi v_0^2}.
\]
В трёх измерениях с тремя ветвями одинаковой скорости
\[
    g_{3\mathrm D}(\omega)=\frac{3V}{(2\pi)^3}\int_0^\infty
       \delta(\omega-v_0k)\,4\pi k^2\,dk
     =\frac{3V}{2\pi^2}\frac{\omega^2}{v_0^3}.
\]
\rresult{При линейном спектре $g(\omega)\propto\omega^{d-1}$.}

\subsection*{Точный спектр одномерной цепочки}

Для цепочки
\[
    \omega(k)=\omega_0\left|\sin\frac{ka}2\right|,\qquad
    \omega_0=2\sqrt{\frac\alpha m},\qquad
    -\frac\pi a<k\le\frac\pi a.
\]

\begin{rbox}[left=18pt,right=18pt,top=12pt,bottom=10pt]{[]}
\centering
\begin{tikzpicture}[>=Latex,x=1.25cm,y=1.6cm,
    every node/.style={text=rtext}]
    \draw[->] (-3.5,0)--(3.6,0) node[right] {$k$};
    \draw[->] (0,-0.1)--(0,1.45) node[above] {$\omega$};
    \draw[rc3,thick,domain=-3.1416:3.1416,samples=81]
        plot (\x,{abs(sin(\x/2 r))});
    \draw[rconceptcolor,dashed] (-3.3,0.6)--(3.3,0.6)
        node[right] {$\omega$};
    \draw[densely dotted] (-1.287,0)--(-1.287,0.6)
        (1.287,0)--(1.287,0.6);
    \fill (-1.287,0.6) circle (1.6pt) (1.287,0.6) circle (1.6pt);
    \node[below] at (-1.287,0) {$-k_0$};
    \node[below] at (1.287,0) {$k_0$};
    \draw[densely dotted] (-3.1416,0)--(-3.1416,1) (3.1416,0)--(3.1416,1);
    \node[below] at (-3.1416,0) {$-\pi/a$};
    \node[below] at (3.1416,0) {$\pi/a$};
    \draw[densely dotted] (-3.1416,1)--(3.1416,1);
    \node[above left] at (0,1) {$\omega_0$};
\end{tikzpicture}
\captionof{figure}{Закон дисперсии цепочки. Частоте $\omega<\omega_0$
отвечают два корня $\pm k_0$.}
\end{rbox}

Используем свойство дельта-функции
\[
    \delta(f(x))=\sum_n\frac{\delta(x-x_n)}{|f'(x_n)|},\qquad f(x_n)=0.
\]
Корни $\pm k_0$ находятся из $\sin(k_0a/2)=\omega/\omega_0$, и в них
\[
    |f'(k_0)|=\frac{\omega_0a}2\cos\frac{k_0a}2
       =\frac a2\sqrt{\omega_0^2-\omega^2}.
\]
Два корня дают одинаковый вклад:
\[
    g_{1\mathrm D}(\omega)=\frac{2L}{2\pi}\int_0^{\pi/a}
       \delta\left(\omega-\omega_0\sin\frac{ka}2\right)dk
     =\frac{2L}{\pi a\sqrt{\omega_0^2-\omega^2}},
     \qquad 0<\omega<\omega_0.
\]
Проверка нормировки: $\int_0^{\omega_0}g\,d\omega
=\frac{2L}{\pi a}\cdot\frac\pi2=\frac La=N$.
При $\omega\ll\omega_0$ получается $g\to2L/(\pi a\omega_0)=L/(\pi v_0)$
с $v_0=\omega_0a/2$, как и для линейного спектра.
На краю зоны $g$ имеет интегрируемую особенность
\rconcept{ван Хова}: там групповая скорость обращается в нуль.

\subsection*{Энергия цепочки}

Подставляя плотность состояний и вводя $z=\hbar\omega/T$,
$z_0=\hbar\omega_0/T$:
\[
    E=\int_0^{\omega_0}\frac{\hbar\omega}{e^{\hbar\omega/T}-1}\,
       \frac{2L\,d\omega}{\pi a\sqrt{\omega_0^2-\omega^2}}
     =\frac{2LT}{\pi a}\int_0^{z_0}
       \frac{z\,dz}{\sqrt{z_0^2-z^2}\,(e^z-1)}.
\]
Здесь опущена не зависящая от температуры нулевая энергия.
Анализ этого интеграла --- на следующем занятии.

\end{document}
